%% =====================================================================
%%  Integrity-Gated Dual-Channel Trust Fusion (IG-DCTF)
%%  Manuscript draft for IEEE Transactions on Information Forensics and Security
%%  Class: IEEEtran (journal mode, two columns).  Build: pdfLaTeX + BibTeX.
%%
%%  Figures are compiled separately (paper/figures/figNN_*.tex -> figNN_*.pdf)
%%  and placed next to this file in paper/figures/. Until a PDF exists the
%%  manuscript shows a labelled placeholder box instead of failing.
%%  All numbers come from paper/results.md (generated by paper/scripts/compute_results.py).
%% =====================================================================
\documentclass[journal]{IEEEtran}

\usepackage[T1]{fontenc}
\usepackage{amsmath,amssymb,amsfonts}
\usepackage{graphicx}
\usepackage{adjustbox}
\usepackage{booktabs}
\usepackage{array}
\usepackage{cite}
\usepackage{algorithm}
\usepackage{algpseudocode}
\usepackage{xcolor}
\usepackage[hidelinks]{hyperref}

\hyphenation{fin-ger-print live-ness}

% -- figure helper: include figures/<name>.pdf at (at most) the given width; placeholder if missing
\newcommand{\figinc}[2]{%
  \IfFileExists{figures/#1.pdf}%
    {\adjustbox{max width=#2}{\includegraphics{figures/#1.pdf}}}%
    {\fbox{\parbox[c][6em][c]{0.9\dimexpr#2\relax}{\centering\footnotesize
      \texttt{\detokenize{#1}.pdf} not found.\\Compile \texttt{\detokenize{figures/#1.tex}}.}}}}

% -- notation
\newcommand{\Tcve}{T_{\mathrm{cve}}}
\newcommand{\Tfin}{T_{\mathrm{final}}}
\newcommand{\thlo}{\theta_{\mathrm{low}}}
\newcommand{\thhi}{\theta_{\mathrm{high}}}
\newcommand{\Drift}{\mathrm{Drift}}
\newcommand{\ind}[1]{\mathbf{1}[#1]}

\begin{document}

\title{Integrity-Gated Dual-Channel Trust Fusion for\\
OSINT-Discovered Surveillance Cameras}

\author{First~Author,~\IEEEmembership{Student Member,~IEEE,}
        Second~Author,
        and~Third~Author% <-- replace with the real author list and affiliations
\thanks{Manuscript received XX XX, 2026; revised XX XX, 2026. (Corresponding author: First Author.)}
\thanks{The authors are with the Department of XXX, Institute XXX, City, Country (e-mail: xxx@xxx).}}

\markboth{IEEE Transactions on Information Forensics and Security,~Vol.~XX, No.~X, Month~2026}%
{Author \MakeLowercase{\textit{et al.}}: Integrity-Gated Dual-Channel Trust Fusion}

\maketitle

%% ---------------------------------------------------------------------
\begin{abstract}
Camera-based alerting systems implicitly treat every incoming feed as authentic, while IoT trust-scoring schemes rate a camera only by its vulnerability posture. A camera with a pristine CVE profile can nevertheless be swapped for a decoy, fed a looped stream, or relayed through a downscaling intermediary, and a vulnerability-only score will still rate it fully trustworthy. We present Integrity-Gated Dual-Channel Trust Fusion (IG-DCTF), which scores a device along two independent channels: a cyber-posture channel (authentication, CVE class, ownership, patch age, corroboration) and a signal-authenticity channel. The latter has two tiers: a passive Tier~1 that measures drift between consecutive OSINT banner fingerprints and needs no additional connection to the camera, and a Tier~2 that tracks perceptual-hash freezes, sharpness drift and declared-resolution mismatch on authorized footage. The channels are fused by a multiplicative gate, so a clean vulnerability profile cannot compensate for a failed authenticity signal. We prove that the gate caps the final score at 30 whenever signal integrity falls to its lower threshold, and that no single changed banner field can reach that cap. On a 120-scenario labelled corpus, IG-DCTF attains an F\textsubscript{1} of 0.976 (recall 0.952, no false alarms) against 0.811 for the strongest vulnerability-only baseline (exact McNemar $p<10^{-4}$), with identical behaviour on the 90 scenarios that carry no authenticity evidence. A reachability analysis of the Tier~2 weights showed that an initial choice let only freezes reach low trust on their own; with the revised weights each confirmed failure can trip the gate. The corpus is synthetic, and the claims await field data.
\end{abstract}

\begin{IEEEkeywords}
Surveillance security, device trust, signal integrity, tamper detection, open-source intelligence, cyber-physical systems.
\end{IEEEkeywords}

%% ---------------------------------------------------------------------
\section{Introduction}
\IEEEPARstart{V}{ideo} surveillance has become an input to automated decisions: alerts from detectors such as YOLOv8 and ByteTrack are routed to dispatchers, corroborated across neighbouring cameras, and logged as evidence. Every such pipeline rests on an unstated premise, namely that the stream it analyses really comes from the camera it is attributed to. Internet-exposed cameras make that premise fragile. Mirai-class botnets recruit cameras at Internet scale~\cite{antonakakis2017mirai,griffioen2020mirai,famera2025mirai}, supply-chain flaws expose tens of millions of devices~\cite{das2022kalay}, and public scanners such as Shodan~\cite{matherly2015shodan} make unpatched endpoints easy to find.

The security literature has responded with vulnerability scoring: CVSS-based audits of camera models~\cite{bernot2025cctv}, CVE taxonomies~\cite{oliver2025ipcam}, automated VAPT tooling~\cite{auti2025vapt}, and certification-style trust indices~\cite{swami2025sciiot}. The vision literature has produced detectors and trackers that assume a trustworthy feed~\cite{zhang2022bytetrack,luna2018aod,nayak2019loitering,yilmazer2025abandoned}. The two threads meet only rarely. The consequence is a concrete blind spot (Fig.~\ref{fig:blind}): a camera that is authenticated, patched, and free of known CVEs receives the maximum vulnerability-based trust score whether it is the genuine device, a decoy occupying the same address, or a recorded loop. Vulnerability posture and feed authenticity are different properties, and a score computed from the first cannot certify the second.

\begin{figure}[!t]
\centering
\figinc{fig01_blind_spot}{\columnwidth}
\caption{The vulnerability-only blind spot. Both feeds share an identical, clean Channel~A profile ($\Tcve=100$). A vulnerability-only scorer therefore rates both as high trust; IG-DCTF additionally measures signal integrity $s$ and gates the decoy to $\Tfin=30$. Values are illustrative and taken from the corpus scenario eval-121.}
\label{fig:blind}
\end{figure}

This paper makes four contributions.
\begin{enumerate}
\item \textbf{Problem and threat model.} We formalise authenticity-aware device trust for OSINT-discovered cameras, including a passive-measurement boundary under which the platform sends no packets to discovered devices (Section~\ref{sec:problem}).
\item \textbf{IG-DCTF.} We introduce a dual-channel score whose channels are fused by a multiplicative gate, and prove three properties: non-compensation, tolerance of single-field benign updates, and bounded sensitivity (Section~\ref{sec:method}).
\item \textbf{Two authenticity tiers.} Tier~1 turns already-collected banner metadata into a drift signal; Tier~2 combines perceptual-hash freeze detection, EWMA sharpness tracking and declared-spec checking. A reachability analysis of the Tier~2 weights determines which failures can drive the final score to low trust on their own, and motivated the weight vector we use.
\item \textbf{Evaluation with an explicit protocol.} We compare against three vulnerability-only baselines, an ablation, and an additive-fusion control on a fixed corpus with Wilson intervals and a paired exact test, and we state the limits of that corpus (Section~\ref{sec:eval}).
\end{enumerate}
To our knowledge, none of the twenty works in our survey fuses a vulnerability-posture score with a feed-authenticity signal in a single per-alert decision; Table~\ref{tab:related} summarises the comparison.

\section{Related Work}
\label{sec:related}
\textbf{Compromise measurement.} Antonakakis \emph{et al.}~\cite{antonakakis2017mirai} and Griffioen and Doerr~\cite{griffioen2020mirai} measured how cameras and other IoT devices are recruited and re-compromised, the latter reporting turnover on the scale of hours to days; Zhang \emph{et al.}~\cite{zhang2020forensics} and Famera \emph{et al.}~\cite{famera2025mirai} analyse server-side artefacts and post-2019 variants. These works are descriptive or attacker-side; none yields a per-alert decision. We use their time-scale finding only to motivate the decay of stale scans (Section~\ref{sec:impl}).

\textbf{Camera vulnerability assessment.} Bernot \emph{et al.} audit Hikvision, Dahua and Avigilon cameras with CVSS~\cite{bernot2025cctv}; Oliver classifies camera CVEs from the NVD~\cite{oliver2025ipcam,nvd}; a Shodan-based study maps exposed cameras to CVEs~\cite{kth2018cameras}; Auti \emph{et al.} automate scanning and exploit validation~\cite{auti2025vapt}. All measure posture, none feed authenticity, and none score individual alerts.

\textbf{IoT trust and reputation.} Surveys organise trust models into inference, weighted-average, learning and game-theoretic families~\cite{iottrust2023slr,biot2026trust,ferraris2024trust} and note static weights, Sybil and collusion exposure, and the cold-start problem. SCI-IoT~\cite{swami2025sciiot} offers a rigorous certification index with critical gates but evaluates a device at procurement, not an alert at run time. We adopt its gate idea for unauthenticated streams in Channel~A.

\textbf{Video analytics.} YOLO-family detectors with ByteTrack or DeepSORT tracking~\cite{zhang2022bytetrack,jocher2023yolov8}, abandoned-object pipelines~\cite{luna2018aod,yilmazer2025abandoned}, multi-camera loitering with re-identification~\cite{nayak2019loitering}, SMS alerting~\cite{rasal2025suspicious} and epipolar multi-view fusion~\cite{liu2025epipolar} all treat each stream as equally reliable; a recent review notes that such systems simply stop working when cameras fail or are damaged~\cite{yolosurvey2025}. Blur, occlusion and freeze detectors exist as engineering primitives; what is missing is their fusion with device vulnerability posture. Our Tier~2 therefore uses standard primitives (perceptual hashing~\cite{zauner2010phash}, variance of the Laplacian~\cite{pech2000focus}) and claims novelty in the fusion and the passive Tier~1, not in the primitives.

\begin{table}[!t]
\caption{Capability Comparison with Representative Work Classes}
\label{tab:related}
\centering
\footnotesize
\setlength{\tabcolsep}{3pt}
\begin{tabular}{@{}p{2.6cm}cccc@{}}
\toprule
Work class & \begin{tabular}[c]{@{}c@{}}Vuln.\\posture\end{tabular} & \begin{tabular}[c]{@{}c@{}}Feed\\authenticity\end{tabular} & \begin{tabular}[c]{@{}c@{}}Fused\\score\end{tabular} & \begin{tabular}[c]{@{}c@{}}Per-alert\\decision\end{tabular}\\
\midrule
Camera vuln.\ assessment~\cite{bernot2025cctv,oliver2025ipcam,kth2018cameras,auti2025vapt} & Yes & No & No & No\\
IoT trust and certification~\cite{swami2025sciiot,ferraris2024trust,biot2026trust} & Partial & No & No & No\\
Video analytics and tracking~\cite{zhang2022bytetrack,nayak2019loitering,yilmazer2025abandoned,rasal2025suspicious} & No & No & No & Yes\\
\textbf{IG-DCTF (this work)} & Yes & Yes & Yes & Yes\\
\bottomrule
\end{tabular}
\end{table}

%% ---------------------------------------------------------------------
\section{Problem Formulation and Threat Model}
\label{sec:problem}
\subsection{Setting}
A platform maintains a set of devices $\mathcal{D}$. For each $d\in\mathcal{D}$ it holds an attribute record $a_d$ (authentication state, CVE identifiers and classes, owner type, last patch date), a time series of banner fingerprints $\mathcal{B}_d(t)$ extracted from OSINT scans, and, for a small authorized subset $\mathcal{D}_A\subseteq\mathcal{D}$, access to recorded or licensed footage. A detection event $e$ raised for device $d$ must be mapped to a trust score $\Tfin(e)\in[0,100]$ and a tier: \emph{high} ($\ge 80$), \emph{medium} ($50$--$79$) or \emph{low} ($<50$). Low-trust events are logged but do not trigger emergency dispatch. We call an event \emph{fabricated} if it originates from a device whose posture or feed should not be relied on, and \emph{genuine} otherwise; \emph{fabricated} is the positive class throughout.

\subsection{Adversaries}
Fig.~\ref{fig:threat} lists four adversary capabilities.
$\mathcal{A}_1$ substitutes the endpoint with a decoy or honeypot, or inserts a relay that re-encodes the stream, while keeping the IP address. $\mathcal{A}_2$ replaces the live stream with a recorded loop or a frozen frame. $\mathcal{A}_3$ controls several adjacent devices and forges mutual corroboration. $\mathcal{A}_4$ submits falsified device attributes in the event payload.

\begin{figure}[!t]
\centering
\figinc{fig03_threat_model}{\columnwidth}
\caption{Threat model: adversary capability, the entry point it uses, and the component that responds. Only $\mathcal{A}_1$ and $\mathcal{A}_2$ are evaluated quantitatively; $\mathcal{A}_3$ and $\mathcal{A}_4$ are covered by functional tests.}
\label{fig:threat}
\end{figure}

\subsection{Assumptions and Boundary}
We assume the platform, its database and its API key are uncompromised, and that the first banner fingerprint recorded for a device is genuine. The platform is \emph{passive}: it only reads metadata that a commercial scanner has already collected and never connects to a discovered device. Frame-level analysis is restricted by construction to a whitelist of authorized streams; a non-whitelisted RTSP source raises an error before any frame is read. An adversary who reproduces the baseline banner exactly (static impersonation), or who degrades a stream slowly enough for the sharpness baseline to follow, is outside what the signals can detect; we return to this in Section~\ref{sec:discussion}.

%% ---------------------------------------------------------------------
\section{Integrity-Gated Dual-Channel Trust Fusion}
\label{sec:method}
\begin{figure*}[!t]
\centering
\figinc{fig02_architecture}{\textwidth}
\caption{IG-DCTF architecture. Channel~A (grey) scores cyber posture from NVD, KEV and ownership data. Channel~B (accent) scores signal authenticity: Tier~1 runs on every device from passive banners, Tier~2 only on authorized footage and receives the declared resolution from Tier~1. The multiplicative gate combines both into $\Tfin$, which drives dispatch, the audit ledger and the dashboard.}
\label{fig:arch}
\end{figure*}

\subsection{Channel A: Cyber Posture}
With $k$ corroborating neighbours, the cyber-posture score is
\begin{equation}
\Tcve=\operatorname{clip}_{[0,U]}\!\Big(100-30\,\ind{\neg\mathrm{auth}}-\Delta_{\mathrm{cat}}-20\,\ind{\mathrm{owner\;unknown}}-15\,\ind{\mathrm{stale}}-10\,\ind{\mathrm{lat}>500}+\beta(k)\Big),
\label{eq:tcve}
\end{equation}
where $U=49$ for an unauthenticated stream and $U=100$ otherwise (a critical gate in the spirit of~\cite{swami2025sciiot}); $\mathrm{stale}$ means the last patch is older than two years or unknown; $\Delta_{\mathrm{cat}}=0$ without CVEs and otherwise $30$ for authentication-bypass or remote-code-execution classes, $25$ for command-injection or memory-corruption classes or two or more CVEs, and $15$ for the rest; and $\beta(k)=-10,\,0,\,+20$ for $k=0,\,1,\,\ge 2$. Device attributes are read from the platform database only, never from the event payload.

\subsection{Channel B, Tier 1: Passive Banner Drift}
\begin{figure}[!t]
\centering
\figinc{fig04_tier1_drift}{\columnwidth}
\caption{Tier~1 worked example. Fields that changed between consecutive scans are shown in the accent colour. Each field contributes $w_i\delta_i$; four changed fields give $\Drift=0.85$, so $s=0.15$ and the gate trips. Values are illustrative.}
\label{fig:tier1}
\end{figure}
A fingerprint is the tuple $\mathcal{B}_d(t)=\langle\text{product},\text{server},\text{firmware},\mathcal{P},\text{resolution},\text{codec}\rangle$ parsed from the banner, where $\mathcal{P}$ is the open-port set. Between consecutive scans,
\begin{equation}
\Drift(d,t)=\frac{\sum_i w_i\,\delta_i\big(\mathcal{B}_d(t)[i],\mathcal{B}_d(t{-}1)[i]\big)}{\sum_i w_i}\in[0,1],
\label{eq:drift}
\end{equation}
with weights $w=(0.10,0.20,0.25,0.20,0.20,0.05)$ for product, server, firmware, ports, resolution and codec. For categorical fields $\delta=0$ if the values agree or both are unknown, $0.5$ if exactly one is unknown, and $1$ otherwise. For ports $\delta=1-|\mathcal{P}_t\cap\mathcal{P}_{t-1}|/|\mathcal{P}_t\cup\mathcal{P}_{t-1}|$. For resolution $\delta=1$ for a downgrade, $0.6$ for any other change, and $0$ otherwise. Extraction and comparison are pure functions of stored metadata (Fig.~\ref{fig:tier1}).

\subsection{Channel B, Tier 2: Frame Liveness}
\begin{figure*}[!t]
\centering
\figinc{fig05_tier2_liveness}{\textwidth}
\caption{Tier~2 liveness tracker, run per authorized stream. A freeze counter over perceptual-hash distances, an EWMA sharpness baseline, and a declared-resolution check each yield a signal in $[0,1]$; the aggregator takes the largest weighted term, with weights $(1.0,0.6,0.5)$.}
\label{fig:tier2}
\end{figure*}
For each frame $f_t$ the tracker computes a 64-bit perceptual hash~\cite{zauner2010phash} (low-frequency $8{\times}8$ DCT block of a $32{\times}32$ grey image, thresholded at its median) and its Hamming distance $d_t$ to the previous hash. Let $c_t$ be the number of consecutive frames with $d_t\le 1$. Then
\begin{equation}
F_t=\ind{c_t\ge 25},\quad
q_t=\operatorname{clip}\!\Big(1-\frac{v_t}{\mu_t}\Big),\quad
\mu_t=\alpha v_t+(1-\alpha)\mu_{t-1},
\label{eq:tier2a}
\end{equation}
where $v_t$ is the variance of the Laplacian~\cite{pech2000focus}, $\alpha=0.05$, and $q_t=0$ while $\mu_t\le 10$. The mismatch flag $m_t\in\{0,1\}$ is set when the frame is smaller than the resolution declared in the Tier~1 banner (e.g.\ width below 1280 or height below 720 for a declared 1080p). The liveness score is
\begin{equation}
L(t)=1-\max\big(1.0\,F_t,\;0.6\,q_t,\;0.5\,m_t\big)\in[0,1].
\label{eq:liveness}
\end{equation}
The weights are chosen by the reachability argument of Proposition~5: a confirmed freeze removes all liveness, maximal blur leaves $L=0.4$, and a resolution mismatch leaves $L=0.5$, so each failure alone lands at or below $\thlo$. The mismatch weight is the smallest of the three because a genuine camera that serves a lower-resolution sub-stream would also trigger it.

\subsection{Fusion by a Multiplicative Gate}
Signal integrity combines the tiers as
\begin{equation}
s=\begin{cases}(1-\Drift)\,L & \text{if Tier 2 is available},\\ 1-\Drift & \text{otherwise},\end{cases}
\label{eq:s}
\end{equation}
and the gate and final score are
\begin{equation}
G(s)=\begin{cases}
1 & s\ge\thhi,\\
0.3+0.7\,\dfrac{s-\thlo}{\thhi-\thlo} & \thlo\le s<\thhi,\\
0.3 & s<\thlo,
\end{cases}
\qquad
\Tfin=\Tcve\,G(s),
\label{eq:gate}
\end{equation}
with $\thlo=0.50$ and $\thhi=0.85$ (Fig.~\ref{fig:gate}). The implementation additionally applies $\min(30,\cdot)$ when $s<\thlo$; because $\Tcve\le 100$ this is implied by~\eqref{eq:gate} and acts as a safeguard.

\begin{figure}[!t]
\centering
\figinc{fig06_gate_function}{\columnwidth}
\caption{The gate $G(s)$. Below $\thlo$ the multiplier sits at its floor of $0.3$, so $\Tfin\le 30$; between the thresholds it tapers linearly; above $\thhi$ it is transparent. The two marked points are the decoy of Fig.~\ref{fig:tier1} and a single-field firmware update.}
\label{fig:gate}
\end{figure}

\subsection{Properties}
\begin{proposition}[Non-compensation]
If $s\le\thlo$ then $\Tfin\le 30$, whatever the device attributes, so the event is low trust.
\end{proposition}
\begin{IEEEproof}
$G(s)\le 0.3$ for $s\le\thlo$ and $\Tcve\le 100$ by~\eqref{eq:tcve}, hence $\Tfin\le 30<50$.
\end{IEEEproof}

\begin{proposition}[Blind spot of Channel-A-only scorers]
Let $e_1,e_2$ be events with equal Channel-A inputs. Any score that is a function of those inputs alone returns the same tier for both. If their labels differ, it errs on at least one.
\end{proposition}
\noindent This is immediate, but it is the formal statement of Fig.~\ref{fig:blind}: no refinement of the vulnerability weights can remove the error; only an additional signal can.

\begin{proposition}[Tolerance of single-field updates]
With Tier~2 clean, a change in any single banner field gives $\Drift\le 0.25$, hence $s\ge 0.75$ and $G\ge 0.8$. The hard region $s<\thlo$ requires $\Drift>0.5$, which needs at least three changed fields.
\end{proposition}
\begin{IEEEproof}
The largest weight is $0.25$, and the two largest sum to $0.45<0.5$; three fields (e.g.\ $0.25+0.20+0.20$) are needed to exceed $0.5$.
\end{IEEEproof}

\begin{proposition}[Bounded sensitivity]
$G$ is continuous and non-decreasing with Lipschitz constant $0.7/(\thhi-\thlo)=2$; hence $|\Delta\Tfin|\le 2\,\Tcve\,|\Delta s|$.
\end{proposition}

\begin{proposition}[Tier-2 reachability]
With the weights of~\eqref{eq:liveness}, $\Drift=0$ and $\Tcve=100$, a single failure yields $L=s=0.0$ (freeze), $0.4$ (maximal blur) or $0.5$ (resolution mismatch), i.e.\ $\Tfin=30$ in every case. More generally, a failure with weight $w$ at full strength gives low trust ($\Tfin<50$ for $\Tcve=100$) iff $s=1-w<0.6$, i.e.\ iff $w>0.4$. The earlier weight vector $(0.5,0.3,0.2)$ violates this for two of the three failures: it gives $s=0.5,0.7,0.8$ and $\Tfin=30,70,90$, so only freezes reached low trust alone.
\end{proposition}
\begin{IEEEproof}
Substitute into~\eqref{eq:liveness}--\eqref{eq:gate}. $G<0.5$ iff $s<0.5+0.35\cdot(0.2/0.7)=0.6$, and $s=L=1-w$ when $\Drift=0$.
\end{IEEEproof}
The proposition fixes only what a failure of full strength achieves; partial blur tapers continuously through $q_t$. Its price is sensitivity: a mismatch flag alone now suffices to reach low trust, which is correct for a relay that downscales the stream and wrong for a benign sub-stream. We could not measure this trade-off without field data (Section~\ref{sec:discussion}).

\begin{algorithm}[!t]
\caption{IG-DCTF scoring of a detection event}
\label{alg:igdctf}
\begin{algorithmic}[1]
\Require device $d$, event $e$, optional liveness $L$
\State $\Tcve\gets$ Eq.~\eqref{eq:tcve} from $a_d$ and $k$ corroborating cameras
\State $\mathcal{B}_{\mathrm{new}}\gets$ parse latest stored banner of $d$; $\mathcal{B}_{\mathrm{old}}\gets$ stored fingerprint
\State $\Drift\gets$ Eq.~\eqref{eq:drift}$(\mathcal{B}_{\mathrm{old}},\mathcal{B}_{\mathrm{new}})$; store $\mathcal{B}_{\mathrm{new}}$
\State $s\gets(1-\Drift)\cdot L$ if $L$ given, else $1-\Drift$
\State $\Tfin\gets\Tcve\cdot G(s)$; if $s<\thlo$ then $\Tfin\gets\min(30,\Tfin)$
\State tier $\gets$ high if $\Tfin\ge 80$, medium if $\ge 50$, else low
\State dispatch by tier; append decision (and a trip entry if $s<\thlo$) to the ledger
\State \Return $\Tfin$, tier, $s$
\end{algorithmic}
\end{algorithm}

%% ---------------------------------------------------------------------
\section{Implementation}
\label{sec:impl}
\begin{figure*}[!t]
\centering
\figinc{fig07_event_pipeline}{\textwidth}
\caption{Request path of a detection event. Stage (a) rejects unauthenticated, replayed, over-rate or stale requests before any scoring; stage (b) computes the score, dispatches by tier, records the decision and broadcasts it. Gate trips additionally produce their own ledger entry and a dedicated dashboard event.}
\label{fig:pipeline}
\end{figure*}
The platform is a FastAPI backend with SQLite (WAL), a React dashboard and a separate video process (YOLOv8 with ByteTrack, rule engine for loitering and perimeter breach). Devices are discovered by passive Shodan queries for three Indian cities, enriched with NVD CVEs, CWE-derived classes, CISA KEV~\cite{cisakev} membership and WHOIS ownership, and cached for 24~h. The request path is shown in Fig.~\ref{fig:pipeline}. Admission uses a constant-time API-key comparison, a five-minute idempotency cache (HTTP 409), a per-camera limit of 10 events per minute (429), and a 60\,s timestamp window (400). Each decision is appended to a SHA-256 hash chain, $H_i=\mathrm{SHA256}(H_{i-1}\,\|\,\mathrm{JSON}(p_i))$ with sorted keys, in the style of a Merkle chain~\cite{merkle1987}; the head hash can be published for external anchoring. Gate trips receive a dedicated ledger entry type.

Three auxiliary modules are implemented and unit-tested but not evaluated here. (i) \emph{Adaptive decay}: scores erode with half-life $T_{1/2}=48\,\mathrm{h}/(1+1.5\,\mathrm{EPSS}+\mathrm{KEV})$, floored at 6\,h, using EPSS~\cite{jacobs2021epss} and the KEV active fraction. (ii) \emph{Cold start}: a device with no history takes its Bayesian prior from a cluster of same-manufacturer, port-similar devices, clamped to $[0.2,0.8]$. (iii) \emph{Collusion analysis}: a velocity counter flags camera pairs that corroborate each other more than five times per hour, and a graph heuristic flags components of four or more cameras whose edge count implies a cycle; an optional two-layer graph convolutional network~\cite{kipf2017gcn} trained on synthetic graphs scores nodes.

%% ---------------------------------------------------------------------
\section{Evaluation}
\label{sec:eval}
Every number below is produced by a single script (\texttt{paper/scripts/compute\_results.py}) and recorded in \texttt{paper/results.md}.

\subsection{Questions}
RQ1: does adding Channel~B reduce missed untrustworthy events relative to vulnerability-only scorers? RQ2: what do the channels and the fusion rule each contribute? RQ3: does the gate degrade scenarios that carry no authenticity evidence? RQ4: how do the two tiers behave in isolation and end to end? RQ5: how sensitive is the result to the thresholds, and what does it cost?

\subsection{Corpus, Baselines and Protocol}
The corpus holds 120 directly scoreable scenarios (63 fabricated, 57 genuine); ten API-layer attack scenarios (replay, stale timestamp, rate flood) are excluded because they test the request path (HTTP 400/409/429) rather than the scorers. Ninety scenarios describe device attributes and corroboration only; thirty also carry a drift value and/or a liveness value (Table~\ref{tab:corpus}). Scenarios and labels were authored by the project team; for the 90 attribute-only scenarios the label follows the intended tier of the scenario, so those scenarios test consistency of the scorers more than they test them against field data. The corpus is therefore a controlled benchmark, not a field study.

Baselines are the weighted-average scorer (WA; fixed deductions of $30,25,20,15,10$ and a $+20$ corroboration bonus), the category-aware Channel-A scorer of~\eqref{eq:tcve} (Advanced), and a Bayesian log-odds scorer with fixed likelihood ratios. A scenario is predicted positive iff its final score is below 50. We report accuracy, recall, specificity and $F_1$, 95\% Wilson intervals~\cite{wilson1927}, and an exact McNemar test~\cite{mcnemar1947} on paired correctness.

\begin{table}[!t]
\caption{Composition of the 120-Scenario Corpus}
\label{tab:corpus}
\centering
\footnotesize
\begin{tabular}{@{}lccc@{}}
\toprule
Subset & Fabricated & Genuine & Total\\
\midrule
Attributes only (no Channel-B field) & 45 & 45 & 90\\
With drift and/or liveness & 18 & 12 & 30\\
\midrule
Total & 63 & 57 & 120\\
\bottomrule
\end{tabular}
\end{table}

\subsection{Main Comparison (RQ1)}
\begin{table*}[!t]
\caption{Four-Model Comparison on the 120 Scenarios (Positive = Fabricated; 95\% Wilson Intervals in Brackets)}
\label{tab:main}
\centering
\footnotesize
\begin{tabular}{@{}lcccccrrrr@{}}
\toprule
Model & Accuracy & Recall & Specificity & $F_1$ & Precision & TP & FP & TN & FN\\
\midrule
WA (baseline)           & 0.792 [0.711, 0.855] & 0.603 [0.480, 0.715] & 1.000 & 0.753 & 1.000 & 38 & 0 & 57 & 25\\
Advanced (Channel A)    & 0.833 [0.757, 0.889] & 0.683 [0.560, 0.784] & 1.000 & 0.811 & 1.000 & 43 & 0 & 57 & 20\\
Bayesian log-odds       & 0.842 [0.766, 0.896] & 0.746 [0.627, 0.837] & 0.947 & 0.832 & 0.940 & 47 & 3 & 54 & 16\\
\textbf{IG-DCTF}        & \textbf{0.975} [0.929, 0.991] & \textbf{0.952} [0.869, 0.984] & 1.000 & \textbf{0.976} & 1.000 & 60 & 0 & 57 & 3\\
\bottomrule
\end{tabular}
\end{table*}

\begin{figure}[!t]
\centering
\figinc{fig08_comparison}{\columnwidth}
\caption{Metric comparison of the four scorers on the 120-scenario corpus. Whiskers are 95\% Wilson intervals (not shown for $F_1$).}
\label{fig:comparison}
\end{figure}

Table~\ref{tab:main} and Fig.~\ref{fig:comparison} show that IG-DCTF misses 3 of 63 untrustworthy events, against 20 for the best vulnerability-only scorer, and raises no false alarm. Its advantage over Advanced is significant under a paired exact test: IG-DCTF is correct and Advanced wrong on 17 scenarios, the reverse never occurs ($p=1.5\times 10^{-5}$). The three remaining misses (eval-026, eval-030, eval-050) are Channel-A borderline cases (a single XSS CVE on an authenticated corporate camera; an unknown-owner camera with one corroborator; an old authenticated camera with two CVEs) that every scorer misclassifies. On this corpus WA and Advanced raise no false alarm, while the Bayesian scorer raises three.

\subsection{Where the Gain Comes From (RQ3)}
\begin{table}[!t]
\caption{Accuracy / $F_1$ Stratified by Evidence Available}
\label{tab:strata}
\centering
\footnotesize
\begin{tabular}{@{}lcc@{}}
\toprule
Model & With Channel-B fields ($n{=}30$) & Without ($n{=}90$)\\
\midrule
WA        & 0.433 / 0.105 & 0.911 / 0.902\\
Advanced  & 0.433 / 0.105 & 0.967 / 0.966\\
Bayesian  & 0.467 / 0.200 & 0.967 / 0.968\\
IG-DCTF   & \textbf{1.000} / \textbf{1.000} & 0.967 / 0.966\\
\bottomrule
\end{tabular}
\end{table}
Table~\ref{tab:strata} separates the scenarios. On the 90 attribute-only scenarios IG-DCTF coincides with Advanced, as it must when $\Drift=0$ and no liveness value is present ($G=1$), so the gate adds no regression there. On the 30 scenarios with authenticity evidence the vulnerability-only scorers are close to chance (they catch 1 of 18 fabricated events), while IG-DCTF classifies all 30 correctly. The aggregate improvement is thus entirely attributable to Channel~B, by construction of the corpus; the experiment demonstrates the blind spot and its closure, not a field accuracy.

\subsection{Ablation and Fusion Rule (RQ2)}
\begin{figure}[!t]
\centering
\figinc{fig09_ablation}{\columnwidth}
\caption{Ablation. Bars show $F_1$; the number in parentheses is the count of missed fabricated events out of 63. No variant raised a false alarm. The additive control subtracts $70(1-s)$ from $\Tcve$, which has the same worst-case end point (30) as the gate.}
\label{fig:ablation}
\end{figure}
Channel~A alone misses 20 events and Channel~B alone (a perfect device scored only through the gate) misses 45 (Fig.~\ref{fig:ablation}); neither substitutes for the other. To test whether multiplicative fusion matters, we also evaluated an additive rule $T=\Tcve-70(1-s)$, whose worst case matches the gate. It misses 10 events (accuracy 0.917, $F_1$ 0.914) against 3 for the gate. The additive rule lets a high $\Tcve$ absorb a moderate integrity failure (a perfect device with $s=0.4$ keeps 58, medium trust, where the gate gives 30), which is exactly what non-compensation forbids.

\subsection{Detector-Level Behaviour (RQ4)}
\begin{table}[!t]
\caption{Detector Benchmarks and End-to-End Consequence ($\Tcve=100$; $\Drift=0$ for the Tier~2 rows)}
\label{tab:detectors}
\centering
\footnotesize
\setlength{\tabcolsep}{3.5pt}
\begin{tabular}{@{}lcccc@{}}
\toprule
Benchmark / category & $n$ & Flagged & Low trust & Mean $\Tfin$\\
\midrule
\multicolumn{5}{@{}l}{\emph{Tier 1 (threshold $\Drift\ge 0.15$; 50 stable, 50 mutated)}}\\
Stable controls                    & 50 & 0  & 0 & 100\\
Single-field mutations (4 types)   & 40 & 40 & 0 & 80--90\\
Full decoy (6 fields)              & 10 & 10 & 10 & 30\\
\addlinespace
\multicolumn{5}{@{}l}{\emph{Tier 2 (80 synthetic sequences)}}\\
Genuine live                       & 30 & 0  & 0  & 100\\
Freeze / loop                      & 20 & 20 & 20 & 30\\
Blur / occlusion                   & 20 & 20 & 20 & 30\\
Resolution mismatch                & 10 & 10 & 10 & 30\\
\bottomrule
\end{tabular}
\end{table}
Table~\ref{tab:detectors} separates \emph{detection} from \emph{consequence}. The Tier~1 detector flags all 50 mutated fingerprints and none of the 50 controls, but the benchmark contains only 13 distinct baseline--current pairs (cases are replicated), and single-field mutations produce $\Drift=0.20$--$0.25$, i.e.\ a taper and not a trip, as Proposition~3 predicts; only the six-field decoy trips the gate. The Tier~2 detector flags all 50 attack sequences and none of 30 controls with a mean latency of 9.6 frames after onset. Pushing the minimum liveness of each run through the gate gives $\Tfin=30$ for all 50 attack sequences and 100 for all 30 controls (mean minimum $L$: 0.00 for freezes, 0.40 for blur, 0.50 for mismatch). This is the intended behaviour of Proposition~5. We report it with a caveat: the weights were revised after a first run with $(0.5,0.3,0.2)$ showed that blur and mismatch sequences ended at $\Tfin=70$ and $90$, so the revision was informed by this very benchmark. The result therefore demonstrates that the design \emph{can} act on each failure, not that the weights generalise. The labelled corpus injects liveness values down to 0.25, which the earlier weights could not produce and the present ones can.

\subsection{Threshold Sensitivity and Cost (RQ5)}
\begin{figure}[!t]
\centering
\figinc{fig10_decision_map}{\columnwidth}
\caption{Decision map in the $(\Drift,L)$ plane. Curves are the iso-lines $(1-\Drift)L=0.50$ and $0.85$; markers are the 30 corpus scenarios that carry Channel-B fields (scenarios without a Tier-2 value are drawn at $L=1$). All genuine scenarios have $s\ge 0.686$ and all fabricated ones $s\le 0.38$.}
\label{fig:map}
\end{figure}
We swept $\thlo\in[0.30,0.60]$ and $\thhi\in[0.70,0.95]$ (42 pairs) with the same tier cut-offs as the main comparison. $F_1$ is 0.976 at every pair. The surface is flat because the corpus has a wide margin around both thresholds (Fig.~\ref{fig:map}); it shows that the result is not tuned to a particular point, but also that this corpus \emph{cannot select} the thresholds. They are justified by design instead: $\thhi=0.85$ lets any single-field update pass with a mild taper (Proposition~3), and $\thlo=0.5$ places the hard region where more than half of the integrity budget is lost. Scenarios near the thresholds need real data. Scoring one event costs 4.7\,$\mu$s for the full pipeline against 2.6\,$\mu$s for Channel~A alone (in-process, excluding I/O).

\subsection{Threats to Validity}
\emph{Construct}: the label ``fabricated'' conflates compromised feeds with weak posture. \emph{Internal}: scenarios and labels were written by the authors; Channel-B values are injected rather than measured, and the Tier~2 weights were revised after inspecting the Tier~2 benchmark (Section~\ref{sec:eval}). \emph{External}: no field data; the Tier~1 benchmark is nearly deterministic; the Tier~2 sequences are synthetic noise, checkerboard and uniform images rather than real video, and day--night illumination changes were not tested. \emph{Statistical}: with 120 scenarios the intervals are wide (Table~\ref{tab:main}), and the McNemar result inherits the corpus design.

%% ---------------------------------------------------------------------
\section{Discussion}
\label{sec:discussion}
\textbf{What the signals cannot see.} Tier~1 compares a device's own banners over time. An adversary who copies the genuine banner exactly (static impersonation), or who controls the banner shown to the scanner, defeats it; Tier~2 is the stronger signal but exists only for authorized streams, which will be a small fraction of discovered cameras. A slowly applied blur is absorbed by the EWMA baseline because the baseline includes the current frame.

\textbf{Weights.} The Tier~2 weights $(1.0,0.6,0.5)$ are set by the reachability condition $w>0.4$ of Proposition~5, not fitted to data, and were revised once after inspecting the synthetic benchmark. Their cost is false trips: a resolution mismatch alone now reaches low trust, so a camera that legitimately serves a sub-stream below its declared resolution would be penalised, and a long static scene could mimic a freeze. Neither effect is measured here. Learning the weights, or making the mismatch term depend on the stream profile, from labelled field data is a further step.

\textbf{Ethics and dual use.} The platform handles infrastructure metadata only, extracts no biometric or personal data, sends no packets to discovered devices, and processes video only from a whitelist. Aggregated vulnerability maps can be misused as target lists; the design mitigates this with authenticated write endpoints, rate limits and a tamper-evident ledger, and critical findings are to be reported to the relevant CERT before dissemination.

\textbf{Future work.} Longitudinal Shodan snapshots with known device replacements would turn Tier~1 from a synthetic to an empirical result; real footage with day--night transitions would test the freeze rule; adaptive adversaries who optimise against the gate should be studied; and the auxiliary modules of Section~\ref{sec:impl} need their own evaluation.

%% ---------------------------------------------------------------------
\section{Conclusion}
Vulnerability posture and feed authenticity are different properties, and a trust score built on the first cannot certify the second. IG-DCTF measures both and fuses them multiplicatively, so that a failed authenticity signal bounds the final score regardless of how clean the CVE profile is. On a controlled corpus this removes the blind spot without side effects on scenarios lacking authenticity evidence. The analysis also shows what remains open: the Tier~2 weights are set by a reachability condition and trade missed attacks against false trips that we could not measure, thresholds cannot be tuned on synthetic data, and the strongest claims await evaluation on field data.

%% ---------------------------------------------------------------------
\bibliographystyle{IEEEtran}
\bibliography{references}

\end{document}
